\chapter*{Zoznam príloh}
\addcontentsline{toc}{chapter}{Zoznam príloh}

\begin{description}[leftmargin=3cm, labelwidth=3cm, labelsep=1em, font=\bfseries]

    \item[Príloha A] Evidencia využitia generatívnej umelej inteligencie

    \item[Príloha B] Zdrojový kód systému

    \item[Príloha C] Fotodokumentácia a~videozáznam inštalácie

\end{description}

% ============================================================================

\appendix

\chapter{Evidencia využitia generatívnej umelej inteligencie}
\label{priloha:ai-evidencia}

Pri spracovaní práce boli využité nástroje generatívnej umelej inteligencie
ako podporné pomôcky. Úplná história použitých promptov nebola systematicky
archivovaná, preto je evidencia spracovaná spätne podľa účelu použitia.
Použité vstupy sa týkali najmä jazykového preformulovania vlastných
podkladov bez dopĺňania nových odborných faktov, zhodnotenia
zrozumiteľnosti kapitol, orientačnej práce s~citáciami a~zdrojmi,
návrhu niektorých implementačných postupov a~pomoci pri debugovaní.

\begin{table}[H]
    \centering
    \small
    \begin{tabular}{|p{2.7cm}|p{2.0cm}|p{5.0cm}|p{3.9cm}|}
        \hline
        \textbf{Nástroj} & \textbf{Obdobie} & \textbf{Účel a~spôsob použitia} & \textbf{Overenie výstupov} \\
        \hline
        ChatGPT (OpenAI) & január--máj 2026 &
        Práca s~textom, jazyková a~štylistická úprava vlastných podkladov,
        lepšia formulácia myšlienok, zhodnotenie zrozumiteľnosti kapitol
        a~orientačná práca s~citáciami. &
        Kontrola, či upravené formulácie nemenia pôvodnú myšlienku,
        porovnanie s~citovanými zdrojmi a~finálna úprava textu. \\
        \hline
        Claude (Anthropic) & 2025--2026 &
        Podpora pri práci so zdrojovým kódom, písaní testov, hľadaní chýb,
        dopĺňaní komentárov a~formulovaní textov logovacích správ. &
        Ručná kontrola upraveného kódu, jeho začlenenie do existujúcej
        štruktúry projektu a~následné testovanie funkčnosti. \\
        \hline
    \end{tabular}
    \caption{Evidencia použitia nástrojov generatívnej umelej inteligencie}
    \label{tab:ai-evidencia}
\end{table}

Nástroje generatívnej umelej inteligencie neboli použité ako samostatný
zdroj odborných tvrdení. Výber, interpretácia a~použitie odborných zdrojov,
návrh systému, hardvérová realizácia, integrácia, testovanie a~finálna
podoba práce vychádzali z~vlastnej realizácie systému.

% ============================================================================

\chapter{Zdrojový kód systému}
\label{priloha:zdrojovy-kod}

Zdrojový kód je priložený na elektronickom médiu v~adresári
\texttt{Zdrojovy\_kod/museum-system/} s~nasledovnou štruktúrou:

\begin{description}[leftmargin=4cm, labelwidth=4cm, labelsep=1em]
    \item[\texttt{esp32/}] Firmvér pre ESP32 uzly (tlačidlový uzol,
        motorový regulátor, relé regulátor) --- Arduino/C++
    \item[\texttt{frontend/}] Webový dashboard pre operátorov múzea --- React, Vite
    \item[\texttt{raspberry\_pi/}] Backend riadiaceho systému MuseumController --- Python~3
\end{description}

% ============================================================================

\chapter{Fotodokumentácia a~videozáznam inštalácie}
\label{priloha:foto}

Fotodokumentácia reálnej inštalácie a~videozáznam overenia synchronizácie
multimédií s~fyzickými efektmi sú priložené na elektronickom médiu.
Fotografie fyzického exponátu sú uložené v~adresári
\texttt{Obrazky/Fyzicka\_instalacia/}, screenshoty webového rozhrania
v~adresári \texttt{Obrazky/Web\_UI/} a~videozáznam v~adresári
\texttt{Video/}.
